% !TeX root = ../../Book2.tex
\section{Gauss 函数与自对偶性}
\label{b2:sec:0602}

Poisson 求和把一个函数的格点和转成其变换的格点和。Gauss 函数的变换仍为 Gauss 函数，因此可以显式比较大尺度和小尺度，并把这种比较放回热传播。

\subsection{Gauss 函数的 Fourier 变换}
\label{b2:subsec:060201}

记 \(g_t(x)=e^{-\pi t x^2}\)，其中 \(t>0\)。由\cref{b2:thm:03-gaussian}及伸缩公式，
\[
 \widehat {g_t}(\xi)=t^{-1/2}e^{-\pi\xi^2/t}.
\]
与保持单位质量的 \(G_\varepsilon\) 相比，\(g_t\) 在原点的高度固定为 \(1\)，总积分为 \(t^{-1/2}\)。将两类参数混用会漏掉平方根因子；这里分别用 \(t\) 和 \(\varepsilon\) 保留区别。

\ProseParagraphBreak

平移和调制给出更一般的波包公式。对实数 \(a,b\)，
\[
 \mathcal F\bigl(e^{-\pi t(x-a)^2}e^{2\pi\mathrm{i}bx}\bigr)(\xi)
 =t^{-1/2}e^{-2\pi\mathrm{i}a(\xi-b)}e^{-\pi(\xi-b)^2/t}.
\]
因此空间中心 \(a\) 产生相位，调制频率 \(b\) 产生频率中心，而宽度参数倒置。每一项都来自\chapref{b2:ch:03}已经证明的运算规则。

\subsection{自对偶结构}
\label{b2:subsec:060202}

若 \(\widehat f=f\)，称 \(f\) 对当前 Fourier 变换\term[zi-duiou]{自对偶}[self-dual]。Gauss 函数 \(g_1\) 满足这一条件。一般 \(g_t\) 的参数变为 \(1/t\)，并附带 \(t^{-1/2}\)；只有 \(t=1\) 保持同一个归一化函数。

\ProseParagraphBreak

自对偶不只出现在正 Gauss 函数上。由 Hermite 特征值公式，所有 \(h_{4j}\) 都自对偶，有限线性组合仍然自对偶。相反，\(h_{4j+2}\) 变换后变号，奇数指标对应虚特征值。自对偶是一项线性变换关系，不意味着非负性或概率密度性质。

\subsection{热核解释}
\label{b2:subsec:060203}

\chapref{b2:ch:05}已使用 \(H_t(x)=(4\pi t)^{-1/2}e^{-x^2/(4t)}\)，其变换为 \(e^{-4\pi^2t\xi^2}\)。在实线上周期化，得到
\begin{equation}\label{b2:eq:06-periodic-heat}
 K_t(x)\coloneqq\sum_{n\in\mathbb Z}H_t(x+n)
 =\sum_{k\in\mathbb Z}e^{-4\pi^2tk^2}e^{2\pi\mathrm{i}kx}.
\end{equation}
第二个表达由 Poisson 求和得到；对固定 \(t>0\)，两种表达都绝对收敛。函数 \(K_t\) 称为圆周上的\term[rehe]{热核}[heat kernel]。
\symidx{\(K_t\)：圆周热核}

\begin{proposition}\label{b2:prop:06-heat-kernel}
圆周热核非负、积分为 \(1\)，满足 \(K_t*K_s=K_{t+s}\) 和 \(\partial_tK_t=\partial_x^2K_t\)。对 \(f\in L^p(\mathbb T)\)、\(1\leq p<\infty\)，有 \(K_t*f\to f\) 于 \(L^p\)；对连续周期函数则一致收敛。
\end{proposition}
\begin{proof}
非负性由实线表达给出，积分为 \(1\) 由求和换序后实线 Gauss 积分给出。频率表达的系数相乘得到半群关系，再用周期系数唯一性。对 \(t\geq\varepsilon>0\)，任意阶 \(x,t\) 导数只在系数上增加一个固定次幂的 \(k\)，仍被可和的指数尾统一控制，故可逐项求导，并得到热方程。

为证初值逼近，对 \(0<\delta<1/2\)，圆周上离原点至少为 \(\delta\) 的区域在周期化前对应到实线中 \(|y|\geq\delta\) 的子集。因此
\[
 \int_{\mathrm{d}_{\mathbb T}(x,0)\geq\delta}K_t(x)\,\mathrm{d}x
 \leq\int_{|y|\geq\delta}H_t(y)\,\mathrm{d}y\longrightarrow0.
\]
所以它是正逼近恒等，\chapref{b2:ch:02}的正核逼近命题给出全部收敛结论。
\end{proof}

小时间时，实线表达用少量邻近格点就能描述局部尖峰；大时间时，频率表达只需保留少量低频，就能看出函数趋向常数 \(1\)。两个公式描述同一个核，却在相反的尺度下各自简洁。下一节的 \(\theta\) 变换正是这种尺度对换在 \(x=0\) 处的表达。
